% ============================================================
%  DOCUMENTO 2: PLAN DE TESIS  (archivo principal)
%  Sistema bidireccional de traduccion de voz Quechua-Espanol con IA preentrenada
%  Compilar:  pdflatex plan_tesis -> biber plan_tesis -> pdflatex plan_tesis -> pdflatex plan_tesis
% ============================================================
\documentclass[12pt,a4paper]{article}
% ============================================================
%  Preambulo compartido - APA7 + espanol + graficos
%  Se incluye con % ============================================================
%  Preambulo compartido - APA7 + espanol + graficos
%  Se incluye con % ============================================================
%  Preambulo compartido - APA7 + espanol + graficos
%  Se incluye con \input{../preambulo.tex} en cada documento.
%  Compilar con: pdflatex -> biber -> pdflatex -> pdflatex
% ============================================================

% ---- Codificacion e idioma ----
\usepackage[utf8]{inputenc}
\usepackage[T1]{fontenc}
\usepackage{lmodern}
\usepackage[spanish,es-tabla]{babel}
\usepackage{csquotes}

% ---- Geometria (margenes tipo tesis; ajustar a SINVIC) ----
\usepackage[a4paper,left=3cm,right=2.5cm,top=2.5cm,bottom=2.5cm]{geometry}

% ---- Interlineado 1.5 ----
\usepackage{setspace}
\onehalfspacing

% ---- Tipografia de titulos ----
\usepackage{titlesec}
\titleformat{\section}{\normalfont\Large\bfseries}{\thesection}{1em}{}
\titleformat{\subsection}{\normalfont\large\bfseries}{\thesubsection}{1em}{}

% ---- Tablas y figuras ----
\usepackage{graphicx}
\usepackage{booktabs}
\usepackage{array}
\usepackage{multirow}
\usepackage{tabularx}
\usepackage{float}
\usepackage[font=small,labelfont=bf]{caption}

% ---- Graficos vectoriales (pgfplots / tikz) ----
\usepackage{pgfplots}
\pgfplotsset{compat=1.18}
\usepackage{tikz}
\usetikzlibrary{shapes.geometric,arrows.meta,positioning,calc}

% ---- Colores ----
\usepackage{xcolor}
\definecolor{azulUNSAAC}{RGB}{0,71,133}
\definecolor{grisSuave}{RGB}{240,240,240}

% ---- Bibliografia APA7 (biblatex + biber) ----
\usepackage[
  backend=biber,
  style=apa,
  sortcites=true,
  sorting=nyt,
  language=spanish
]{biblatex}
\DeclareLanguageMapping{spanish}{spanish-apa}
\addbibresource{../referencias.bib}

% ---- Listas e indices (algoritmos, ecuaciones, anexos; tablas largas) ----
% IMPORTANTE: tocloft debe cargarse ANTES de hyperref.
\usepackage{longtable}
\usepackage{algorithm}
\usepackage{algpseudocode}
\usepackage{tocloft}

% ---- Hipervinculos ----
\usepackage[hidelinks]{hyperref}
\hypersetup{
  colorlinks=true,
  linkcolor=black,
  citecolor=azulUNSAAC,
  urlcolor=azulUNSAAC
}

% ---- Icono ORCID (dibujado con tikz, compatible con pdflatex) ----
\usepackage{orcidlink}

% ---- Encabezado / pie de pagina ----
\usepackage{fancyhdr}
\pagestyle{fancy}
\fancyhf{}
\fancyfoot[C]{\thepage}
\renewcommand{\headrulewidth}{0pt}

% ---- Sangria de primer parrafo ----
\usepackage{indentfirst}
 en cada documento.
%  Compilar con: pdflatex -> biber -> pdflatex -> pdflatex
% ============================================================

% ---- Codificacion e idioma ----
\usepackage[utf8]{inputenc}
\usepackage[T1]{fontenc}
\usepackage{lmodern}
\usepackage[spanish,es-tabla]{babel}
\usepackage{csquotes}

% ---- Geometria (margenes tipo tesis; ajustar a SINVIC) ----
\usepackage[a4paper,left=3cm,right=2.5cm,top=2.5cm,bottom=2.5cm]{geometry}

% ---- Interlineado 1.5 ----
\usepackage{setspace}
\onehalfspacing

% ---- Tipografia de titulos ----
\usepackage{titlesec}
\titleformat{\section}{\normalfont\Large\bfseries}{\thesection}{1em}{}
\titleformat{\subsection}{\normalfont\large\bfseries}{\thesubsection}{1em}{}

% ---- Tablas y figuras ----
\usepackage{graphicx}
\usepackage{booktabs}
\usepackage{array}
\usepackage{multirow}
\usepackage{tabularx}
\usepackage{float}
\usepackage[font=small,labelfont=bf]{caption}

% ---- Graficos vectoriales (pgfplots / tikz) ----
\usepackage{pgfplots}
\pgfplotsset{compat=1.18}
\usepackage{tikz}
\usetikzlibrary{shapes.geometric,arrows.meta,positioning,calc}

% ---- Colores ----
\usepackage{xcolor}
\definecolor{azulUNSAAC}{RGB}{0,71,133}
\definecolor{grisSuave}{RGB}{240,240,240}

% ---- Bibliografia APA7 (biblatex + biber) ----
\usepackage[
  backend=biber,
  style=apa,
  sortcites=true,
  sorting=nyt,
  language=spanish
]{biblatex}
\DeclareLanguageMapping{spanish}{spanish-apa}
\addbibresource{../referencias.bib}

% ---- Listas e indices (algoritmos, ecuaciones, anexos; tablas largas) ----
% IMPORTANTE: tocloft debe cargarse ANTES de hyperref.
\usepackage{longtable}
\usepackage{algorithm}
\usepackage{algpseudocode}
\usepackage{tocloft}

% ---- Hipervinculos ----
\usepackage[hidelinks]{hyperref}
\hypersetup{
  colorlinks=true,
  linkcolor=black,
  citecolor=azulUNSAAC,
  urlcolor=azulUNSAAC
}

% ---- Icono ORCID (dibujado con tikz, compatible con pdflatex) ----
\usepackage{orcidlink}

% ---- Encabezado / pie de pagina ----
\usepackage{fancyhdr}
\pagestyle{fancy}
\fancyhf{}
\fancyfoot[C]{\thepage}
\renewcommand{\headrulewidth}{0pt}

% ---- Sangria de primer parrafo ----
\usepackage{indentfirst}
 en cada documento.
%  Compilar con: pdflatex -> biber -> pdflatex -> pdflatex
% ============================================================

% ---- Codificacion e idioma ----
\usepackage[utf8]{inputenc}
\usepackage[T1]{fontenc}
\usepackage{lmodern}
\usepackage[spanish,es-tabla]{babel}
\usepackage{csquotes}

% ---- Geometria (margenes tipo tesis; ajustar a SINVIC) ----
\usepackage[a4paper,left=3cm,right=2.5cm,top=2.5cm,bottom=2.5cm]{geometry}

% ---- Interlineado 1.5 ----
\usepackage{setspace}
\onehalfspacing

% ---- Tipografia de titulos ----
\usepackage{titlesec}
\titleformat{\section}{\normalfont\Large\bfseries}{\thesection}{1em}{}
\titleformat{\subsection}{\normalfont\large\bfseries}{\thesubsection}{1em}{}

% ---- Tablas y figuras ----
\usepackage{graphicx}
\usepackage{booktabs}
\usepackage{array}
\usepackage{multirow}
\usepackage{tabularx}
\usepackage{float}
\usepackage[font=small,labelfont=bf]{caption}

% ---- Graficos vectoriales (pgfplots / tikz) ----
\usepackage{pgfplots}
\pgfplotsset{compat=1.18}
\usepackage{tikz}
\usetikzlibrary{shapes.geometric,arrows.meta,positioning,calc}

% ---- Colores ----
\usepackage{xcolor}
\definecolor{azulUNSAAC}{RGB}{0,71,133}
\definecolor{grisSuave}{RGB}{240,240,240}

% ---- Bibliografia APA7 (biblatex + biber) ----
\usepackage[
  backend=biber,
  style=apa,
  sortcites=true,
  sorting=nyt,
  language=spanish
]{biblatex}
\DeclareLanguageMapping{spanish}{spanish-apa}
\addbibresource{../referencias.bib}

% ---- Listas e indices (algoritmos, ecuaciones, anexos; tablas largas) ----
% IMPORTANTE: tocloft debe cargarse ANTES de hyperref.
\usepackage{longtable}
\usepackage{algorithm}
\usepackage{algpseudocode}
\usepackage{tocloft}

% ---- Hipervinculos ----
\usepackage[hidelinks]{hyperref}
\hypersetup{
  colorlinks=true,
  linkcolor=black,
  citecolor=azulUNSAAC,
  urlcolor=azulUNSAAC
}

% ---- Icono ORCID (dibujado con tikz, compatible con pdflatex) ----
\usepackage{orcidlink}

% ---- Encabezado / pie de pagina ----
\usepackage{fancyhdr}
\pagestyle{fancy}
\fancyhf{}
\fancyfoot[C]{\thepage}
\renewcommand{\headrulewidth}{0pt}

% ---- Sangria de primer parrafo ----
\usepackage{indentfirst}


% ----- Colores de la "ecuacion universal" del titulo -----
\definecolor{colProceso}{RGB}{0,71,133}    % azul  -> Proceso / Accion
\definecolor{colMetodo}{RGB}{181,128,0}    % ambar -> Tecnologia / Metodo
\definecolor{colContexto}{RGB}{176,42,42}  % rojo  -> Contexto / Aplicacion

% ============================================================
%  INTERRUPTOR DE INDICES PRELIMINARES
%  ---> Una sola instruccion controla si TODOS los indices
%       aparecen (o no) antes de la introduccion:
%         \mostrarindicestrue   = SI se muestran
%         \mostrarindicesfalse  = NO se muestran
% ============================================================
\newif\ifmostrarindices
\mostrarindicestrue

% ----- Nombres en espanol de los indices -----
\renewcommand{\contentsname}{\'Indice general}
\floatname{algorithm}{Algoritmo}
\renewcommand{\listalgorithmname}{\'Indice de algoritmos}

% ----- Indice de ecuaciones (lista propia via tocloft) -----
\newlistof{ecuaciones}{loe}{\'Indice de ecuaciones}
\newcommand{\regecuacion}[1]{%
  \addcontentsline{loe}{ecuaciones}{\protect\numberline{\theequation}#1}}

% ----- Indice de anexos (lista propia via tocloft) -----
\newlistof{anexos}{lanx}{\'Indice de anexos}
\newcommand{\anexo}[1]{%
  \refstepcounter{anexos}%
  \section*{Anexo \Alph{anexos}: #1}%
  \addcontentsline{toc}{section}{Anexo \Alph{anexos}: #1}%
  \addcontentsline{lanx}{anexos}{\protect\numberline{\Alph{anexos}}#1}}

\begin{document}

% ------------------- CARATULA -------------------
\begin{titlepage}
\centering
\singlespacing

% ----- Logo UNSAAC (coloca logo_unsaac.png en la raiz o en esta carpeta) -----
\IfFileExists{../logo_unsaac.png}{\includegraphics[width=2.7cm]{../logo_unsaac.png}}{%
  \IfFileExists{logo_unsaac.png}{\includegraphics[width=2.7cm]{logo_unsaac.png}}{}}
\par\vspace{0.5cm}

{\scshape\bfseries Universidad Nacional de San Antonio Abad del Cusco\par}
\vspace{0.2cm}
\rule{\linewidth}{0.4pt}\par
\vspace{1cm}

{\scshape Plan de tesis\par}
\vspace{1cm}

\rule{\linewidth}{0.4pt}\par
\vspace{0.4cm}
{\large\bfseries%
\textcolor{colProceso}{\MakeUppercase{Sistema bidireccional de traducci{\'o}n
de voz Quechua--Espa{\~n}ol}}
\textcolor{colMetodo}{\MakeUppercase{con IA preentrenada}}
\textcolor{colContexto}{\MakeUppercase{para consultas m{\'e}dicas en Cusco}}\par}
\vspace{0.4cm}
\rule{\linewidth}{0.4pt}\par

\vfill

\begin{flushright}
\textbf{Para optar al t{\'i}tulo profesional de:}\\[2pt]
Ingeniero Inform{\'a}tico y de Sistemas\\[12pt]
\textbf{Presentado por:}\\[2pt]
Alain Armando Condori Contreras~\orcidlink{0000-0000-0000-0000}\\[1pt]
{\small 0000-0000-0000-0000}\\[12pt]% <-- ORCID del autor: reemplazar por el real
\textbf{Asesor:}\\[2pt]
Mtro. Ing. Ra{\'u}l Huillca Huallparimachi~\orcidlink{0000-0000-0000-0000}\\[1pt]
{\small 0000-0000-0000-0000}% <-- ORCID del asesor: reemplazar por el real
\end{flushright}

\vfill

{Per{\'u}, Octubre del 2026\par}
\vspace{0.4cm}
\rule{\linewidth}{0.4pt}\par
\vspace{0.2cm}
{\footnotesize\scshape Escuela Profesional de Ingenier{\'i}a Inform{\'a}tica y de Sistemas\\
Facultad de Inform{\'a}tica\par}
\end{titlepage}

% ------------------- RESUMEN -------------------
\section*{Resumen}
\addcontentsline{toc}{section}{Resumen}
En la regi{\'o}n Cusco, m{\'a}s del 55\,\% de la poblaci{\'o}n tiene el quechua como lengua
materna y cerca de 70\,000 asegurados de EsSalud son quechuahablantes \parencite{essalud2023}.
Esta barrera idiom{\'a}tica afecta la calidad de la atenci{\'o}n en salud \parencite{begazo2026}.
El presente trabajo propone desarrollar y evaluar un sistema bidireccional de traducci{\'o}n
de voz Quechua--Espa{\~n}ol para consultas m{\'e}dicas, que encadena reconocimiento de voz
(MMS), traducci{\'o}n autom{\'a}tica (NLLB adaptado al quechua de Cusco, \texttt{quz}) y
s{\'i}ntesis de voz (MMS-TTS) mediante el ajuste (fine-tuning) de modelos preentrenados de
c{\'o}digo abierto. La salida por voz en ambos sentidos busca incluir a la poblaci{\'o}n
quechuahablante con baja o nula alfabetizaci{\'o}n. La validaci{\'o}n se realizar{\'a} con
consultas simuladas, sin datos reales de pacientes.

\vspace{0.3cm}
\noindent\textbf{Palabras clave:} traducci{\'o}n autom{\'a}tica, quechua, PLN, inclusi{\'o}n digital, comunicaci{\'o}n m{\'e}dico-paciente, ASR, TTS.

% ============================================================
%  PAGINAS PRELIMINARES  ---  controladas por \ifmostrarindices
% ============================================================
\ifmostrarindices

\newpage
\tableofcontents                                   % 1. Indice general

\clearpage
\addcontentsline{toc}{section}{\'Indice de figuras}
\listoffigures                                     % 2. Indice de figuras

\clearpage
\addcontentsline{toc}{section}{\'Indice de tablas}
\listoftables                                      % 3. Indice de tablas

\clearpage
\addcontentsline{toc}{section}{\'Indice de algoritmos}
\listofalgorithms                                  % 4. Indice de algoritmos

\clearpage                                          % 5. Indice de simbolos
\section*{\'Indice de s{\'i}mbolos}
\addcontentsline{toc}{section}{\'Indice de s{\'i}mbolos}
\begin{longtable}{@{}p{2.6cm} p{9cm} p{2.4cm}@{}}
\toprule
\textbf{S{\'i}mbolo} & \textbf{Descripci{\'o}n} & \textbf{Unidad} \\
\midrule \endhead
$\alpha$ & Nivel de significancia estad{\'i}stica & --- \\
$p$      & Valor de probabilidad (p-valor) & --- \\
$N$      & Tama{\~n}o de muestra / n.\ de palabras de referencia & --- \\
$\beta$  & Factor de ponderaci{\'o}n recall/precisi{\'o}n en chrF & --- \\
\bottomrule
\end{longtable}

\clearpage                                          % 6. Indice de abreviaturas
\section*{\'Indice de abreviaturas}
\addcontentsline{toc}{section}{\'Indice de abreviaturas}
\begin{longtable}{@{}p{3cm} p{12cm}@{}}
\toprule
\textbf{Sigla} & \textbf{Significado} \\
\midrule \endhead
IA      & Inteligencia artificial \\
PLN     & Procesamiento del lenguaje natural \\
ASR     & Reconocimiento autom{\'a}tico del habla (\emph{Automatic Speech Recognition}) \\
STT     & Voz a texto (\emph{Speech-to-Text}) \\
TTS     & S{\'i}ntesis de voz (\emph{Text-to-Speech}) \\
MT      & Traducci{\'o}n autom{\'a}tica (\emph{Machine Translation}) \\
NMT     & Traducci{\'o}n autom{\'a}tica neuronal \\
MMS     & \emph{Massively Multilingual Speech} (Meta) \\
NLLB    & \emph{No Language Left Behind} (Meta) \\
WER     & Tasa de error de palabra (\emph{Word Error Rate}) \\
BLEU    & \emph{Bilingual Evaluation Understudy} \\
chrF    & \emph{Character n-gram F-score} \\
MOS     & \emph{Mean Opinion Score} \\
quz     & Quechua Collao / Cusco (c{\'o}digo ISO 639-3) \\
quy     & Quechua Ayacucho (c{\'o}digo ISO 639-3) \\
INEI    & Instituto Nacional de Estad{\'i}stica e Inform{\'a}tica \\
EsSalud & Seguro Social de Salud del Per{\'u} \\
UNSAAC  & Universidad Nacional de San Antonio Abad del Cusco \\
\bottomrule
\end{longtable}

\clearpage
\addcontentsline{toc}{section}{\'Indice de anexos}
\listofanexos                                      % 7. Indice de anexos

\clearpage                                          % 8. Glosario de terminos
\section*{Glosario de t{\'e}rminos}
\addcontentsline{toc}{section}{Glosario de t{\'e}rminos}
\begin{longtable}{@{}p{3.6cm} p{11.4cm}@{}}
\toprule
\textbf{T{\'e}rmino} & \textbf{Definici{\'o}n} \\
\midrule \endhead
Corpus                  & Conjunto estructurado de textos o grabaciones usado para entrenar o evaluar modelos. \\
Fine-tuning             & Ajuste de un modelo preentrenado con datos de un dominio espec{\'i}fico. \\
Lengua de bajo recurso  & Idioma con escasez de datos digitales para entrenar modelos de PLN. \\
Modelo preentrenado     & Modelo entrenado sobre grandes vol{\'u}menes de datos, reutilizable mediante ajuste. \\
Lengua polisint{\'e}tica & Lengua que concentra mucha informaci{\'o}n gramatical en una sola palabra, como el quechua. \\
Pipeline                & Secuencia encadenada de etapas (voz $\to$ texto $\to$ traducci{\'o}n $\to$ voz). \\
\bottomrule
\end{longtable}

\clearpage
\addcontentsline{toc}{section}{\'Indice de ecuaciones}
\listofecuaciones                                  % 9. Indice de ecuaciones

\clearpage
\fi
% ============================================================
%  FIN DE PAGINAS PRELIMINARES
% ============================================================

% ============================================================
\section{Cap{\'i}tulo I: Introducci{\'o}n}

\subsection{Contexto y problema general}
El Per{\'u} es un pa{\'i}s plurilingüe en el que el quechua es la lengua originaria con mayor
n{\'u}mero de hablantes. Los Censos Nacionales 2017 registran a millones de peruanos que
declaran el quechua como lengua materna, con una marcada concentraci{\'o}n en la sierra sur
\parencite{inei2017}. En la regi{\'o}n Cusco esta presencia es particularmente intensa: m{\'a}s
del 55\,\% de la poblaci{\'o}n tiene el quechua como lengua materna \parencite{essalud2023}.
Lejos de constituir un dato meramente demogr{\'a}fico, esta realidad ling{\"u}{\'i}stica
condiciona el acceso efectivo a los servicios p{\'u}blicos esenciales ---la atenci{\'o}n en
salud entre ellos--- cuando tales servicios operan de manera predominante en espa{\~n}ol.

En el {\'a}mbito sanitario, la distancia idiom{\'a}tica entre un personal de salud
mayoritariamente castellanohablante y una poblaci{\'o}n usuaria quechuahablante origina una
barrera comunicativa con consecuencias cl{\'i}nicas directas. EsSalud reconoce que cerca de
70\,000 de sus asegurados en Cusco son quechuahablantes ---alrededor del 15\,\% del total
regional--- y, ante la dificultad de atenderlos adecuadamente, reactiv{\'o} en 2023 la central
telef{\'o}nica en quechua \emph{Kachkaniraqmi} \parencite{essalud2023}. No obstante, esta medida
atiende tr{\'a}mites y orientaci{\'o}n a distancia, y no alcanza el momento m{\'a}s cr{\'i}tico de
la atenci{\'o}n: la interacci{\'o}n cara a cara en el consultorio, donde se recogen los
s{\'i}ntomas, se formula el diagn{\'o}stico y se entregan las indicaciones terap{\'e}uticas.

La evidencia disponible confirma la gravedad del problema. \textcite{begazo2026} documenta de
forma cualitativa c{\'o}mo la desconexi{\'o}n idiom{\'a}tica y cultural vulnera la relaci{\'o}n
terap{\'e}utica en hospitales de referencia de Cusco, con efectos sobre la confianza del
paciente, la adherencia al tratamiento y la seguridad del acto m{\'e}dico. En la misma
direcci{\'o}n, \textcite{cuaila2022} halla una relaci{\'o}n entre el nivel de conocimiento del
quechua del personal asistencial y la calidad de atenci{\'o}n percibida por los pacientes en
{\'A}ncash. Ambos trabajos coinciden en el diagn{\'o}stico ---la barrera ling{\"u}{\'i}stica
deteriora la atenci{\'o}n--- pero se detienen en {\'e}l: ninguno propone ni eval{\'u}a una
herramienta tecnol{\'o}gica que permita mitigarla. Ese es, precisamente, el vac{\'i}o que esta
investigaci{\'o}n aborda.

Las respuestas que hoy se ensayan frente a esta barrera resultan insuficientes o poco
escalables. La disponibilidad de personal asistencial biling{\"u}e es limitada y heterog{\'e}nea
---\textcite{cuaila2022} evidencia que el dominio del quechua var{\'i}a ampliamente entre el
personal y condiciona la calidad percibida de la atenci{\'o}n---, mientras que recurrir a
familiares o acompa{\~n}antes como int{\'e}rpretes improvisados compromete la confidencialidad,
introduce errores de interpretaci{\'o}n y traslada al paciente una carga que no le corresponde.
Las iniciativas institucionales existentes, como la central telef{\'o}nica \emph{Kachkaniraqmi}
\parencite{essalud2023}, operan fuera del consultorio y no asisten el di{\'a}logo cl{\'i}nico en
tiempo real. Se configura as{\'i} la necesidad de una soluci{\'o}n tecnol{\'o}gica espec{\'i}fica,
disponible en el punto de atenci{\'o}n y dise{\~n}ada para el vocabulario propio de la consulta
m{\'e}dica.

A la barrera idiom{\'a}tica se superpone una barrera de alfabetizaci{\'o}n que agrava la
exclusi{\'o}n. Cerca del 22.8\,\% de las mujeres rurales mayores de 15 a{\~n}os no sabe leer ni
escribir \parencite{inei2023}, y solo en la provincia cusque{\~n}a de Quispicanchi se registran
m{\'a}s de 11\,000 personas en condici{\'o}n de analfabetismo \parencite{dispurse2023}. Este
dato tiene una implicancia de dise{\~n}o decisiva: una soluci{\'o}n basada {\'u}nicamente en
texto escrito excluir{\'i}a justamente a una parte de la poblaci{\'o}n que m{\'a}s necesita el
apoyo. Por ello, la salida por voz ---en ambos sentidos del di{\'a}logo--- no es un a{\~n}adido
opcional, sino la condici{\'o}n que hace accesible el sistema a usuarios que no leen
(Figura~\ref{fig:barreras}).

\begin{figure}[H]
\centering
\begin{tikzpicture}
\begin{axis}[
    ybar, width=12cm, height=6.5cm, bar width=28pt,
    ymin=0, ymax=100, ylabel={Porcentaje (\%)},
    symbolic x coords={Lengua materna,Asegurados EsSalud,Analfab. mujer rural},
    xtick=data, nodes near coords, nodes near coords style={font=\small},
    x tick label style={align=center,font=\small,text width=2.5cm},
    enlarge x limits=0.25, ymajorgrids=true, grid style=dashed,
]
\addplot[fill=azulUNSAAC] coordinates {
    (Lengua materna,55)(Asegurados EsSalud,15)(Analfab. mujer rural,22.8)
};
\end{axis}
\end{tikzpicture}
\caption{Magnitud de la problem{\'a}tica en Cusco. Elaboraci{\'o}n propia con datos de \textcite{essalud2023} e \textcite{inei2023}.}
\label{fig:barreras}
\end{figure}

\subsection{Tecnolog{\'i}a o m{\'e}todo}
Paralelamente, el estado del arte en procesamiento del lenguaje natural ofrece hoy una
oportunidad que no exist{\'i}a hace pocos a{\~n}os. La estrategia dominante para lenguas de bajo
recurso ya no consiste en entrenar modelos desde cero, sino en adaptar (fine-tuning) grandes
modelos multiling{\"u}es preentrenados sobre enormes vol{\'u}menes de datos. En la voz, el
proyecto MMS \parencite{pratap2024mms} cubre el quechua de Cusco (\texttt{quz}) tanto en
reconocimiento como en s{\'i}ntesis, lo que habilita una entrada y una salida habladas con un
{\'u}nico ecosistema de modelos; Whisper \parencite{radford2023whisper}, en cambio, declara una
cobertura marginal del quechua y se reserva como l{\'i}nea base. En traducci{\'o}n autom{\'a}tica,
NLLB \parencite{costajussa2022nllb} incluye el quechua, pero en su variante de Ayacucho
(\texttt{quy}) y no la de Cusco, por lo que el par debe adaptarse a \texttt{quz}; esta brecha
dialectal constituye el n{\'u}cleo del aporte t{\'e}cnico de este trabajo. Existen, adem{\'a}s,
antecedentes espec{\'i}ficos de traducci{\'o}n neural para esta lengua polisint{\'e}tica
\parencite{ortega2020quechua,mager2021americasnlp} y recursos de datos peruanos, como el corpus
\emph{Huqariq} \parencite{zevallos2022huqariq}, que habilitan el ajuste de los modelos de voz.

\subsection{Vac{\'i}o de investigaci{\'o}n}
No obstante, persiste una desconexi{\'o}n entre ambos mundos: el problema social est{\'a} bien
documentado y la tecnolog{\'i}a base existe, pero nadie ha unido los dos extremos. No se dispone
de un corpus cl{\'i}nico Quechua--Espa{\~n}ol ni de una evaluaci{\'o}n emp{\'i}rica del
desempe{\~n}o de estos modelos sobre el vocabulario propio de la consulta m{\'e}dica (triaje,
s{\'i}ntomas e indicaciones). En consecuencia, este trabajo se propone desarrollar y evaluar un
sistema bidireccional de traducci{\'o}n de voz Quechua--Espa{\~n}ol para consultas m{\'e}dicas,
orientado al dominio cl{\'i}nico y validado con m{\'e}tricas estandarizadas de reconocimiento y
traducci{\'o}n (WER, BLEU y chrF) \parencite{papineni2002bleu,popovic2015chrf}, como un aporte
concreto a la equidad en el acceso a la salud de la poblaci{\'o}n quechuahablante de Cusco.

\subsection{Problema espec{\'i}fico (formulaci{\'o}n del problema)}
\textbf{Problema general:} {\textquestiondown}De qu{\'e} manera un sistema de traducci{\'o}n
bidireccional Quechua--Espa{\~n}ol por voz, basado en modelos de IA preentrenados, influye en
la comunicaci{\'o}n m{\'e}dico-paciente en consultas hospitalarias de Cusco?

\textbf{Problemas espec{\'i}ficos:}
\begin{enumerate}
  \item {\textquestiondown}Qu{\'e} caracter{\'i}sticas debe reunir un corpus de vocabulario cl{\'i}nico Quechua--Espa{\~n}ol para el ajuste de los modelos de traducci{\'o}n por voz?
  \item {\textquestiondown}En qu{\'e} medida el ajuste (fine-tuning) de modelos preentrenados de reconocimiento y traducci{\'o}n mejora la precisi{\'o}n sobre el vocabulario cl{\'i}nico?
  \item {\textquestiondown}Qu{\'e} modalidad de salida (texto para el personal; texto y voz para el paciente) resulta m{\'a}s confiable y comprensible?
  \item {\textquestiondown}Qu{\'e} nivel de utilidad, confianza y comprensi{\'o}n perciben el personal y los pacientes al usar el sistema en consultas simuladas?
\end{enumerate}

\subsection{Propuesta, objetivo y justificaci{\'o}n}
\textbf{Propuesta.} Se propone desarrollar y evaluar un sistema bidireccional de traducci{\'o}n de
voz Quechua--Espa{\~n}ol para consultas m{\'e}dicas, que encadena reconocimiento de voz (MMS),
traducci{\'o}n autom{\'a}tica (NLLB adaptado de \texttt{quy} a \texttt{quz}) y s{\'i}ntesis de voz
(MMS-TTS) mediante el ajuste de modelos preentrenados de c{\'o}digo abierto, con salida por voz en
ambos sentidos del di{\'a}logo.

\textbf{Objetivo general:} desarrollar y evaluar un sistema de traducci{\'o}n bidireccional
Quechua--Espa{\~n}ol por voz, basado en modelos de IA preentrenados, para mejorar la
comunicaci{\'o}n m{\'e}dico-paciente en consultas hospitalarias de Cusco.

\textbf{Objetivos espec{\'i}ficos:}
\begin{enumerate}
  \item Construir y validar un corpus de vocabulario y frases de consulta m{\'e}dica en Quechua--Espa{\~n}ol.
  \item Adaptar (fine-tuning) los modelos preentrenados de voz (MMS, para reconocimiento y s{\'i}ntesis) y, como aporte central, el modelo de traducci{\'o}n NLLB del quechua de Ayacucho (\texttt{quy}) al de Cusco (\texttt{quz}) sobre el dominio cl{\'i}nico.
  \item Integrar el pipeline bidireccional con salida por voz y texto para ambos interlocutores, de modo que el sistema sea utilizable por pacientes con baja o nula alfabetizaci{\'o}n.
  \item Validar el sistema en consultas simuladas, midiendo precisi{\'o}n, calidad de voz, tiempo y percepci{\'o}n.
\end{enumerate}

\textbf{Justificaci{\'o}n.} \textbf{Te{\'o}rica:} evidencia emp{\'i}rica del desempe{\~n}o de modelos
preentrenados en una lengua de bajo recurso. \textbf{Pr{\'a}ctica:} prototipo piloteable en un
centro de salud. \textbf{Social:} ataca una inequidad documentada \parencite{begazo2026}.
\textbf{Metodol{\'o}gica:} proceso replicable en otras variantes o dominios.

\subsection{Arquitectura propuesta}
La Figura~\ref{fig:arq} muestra el flujo del sistema y el Algoritmo~\ref{alg:pipeline} detalla el
procesamiento de un turno de di{\'a}logo.

\begin{figure}[H]
\centering
\begin{tikzpicture}[
    node distance=0.6cm and 0.7cm,
    box/.style={rectangle, rounded corners, draw=azulUNSAAC, thick, fill=grisSuave, text width=2.0cm, align=center, minimum height=1.2cm, font=\scriptsize},
    core/.style={rectangle, rounded corners, draw=colContexto, very thick, fill=grisSuave, text width=2.0cm, align=center, minimum height=1.2cm, font=\scriptsize},
    io/.style={ellipse, draw=black, thick, fill=white, text width=1.8cm, align=center, font=\scriptsize},
    bi/.style={{Stealth[length=2.5mm]}-{Stealth[length=2.5mm]}, thick}
]
\node[io] (pac) {Paciente\\(voz quechua)};
\node[box, right=of pac] (v1) {MMS\\STT / TTS\\(\texttt{quz})};
\node[core, right=of v1] (mt) {Traducci{\'o}n\\NLLB\\\texttt{quz}$\leftrightarrow$es};
\node[box, right=of mt] (v2) {MMS\\STT / TTS\\(es)};
\node[io, right=of v2] (med) {M{\'e}dico\\(voz espa{\~n}ol)};

\draw[bi] (pac) -- (v1);
\draw[bi] (v1) -- (mt);
\draw[bi] (mt) -- (v2);
\draw[bi] (v2) -- (med);
\end{tikzpicture}
\caption{Pipeline modular \textbf{bidireccional} por voz: entrada y salida habladas en ambos extremos (MMS, STT/TTS) para incluir a usuarios con baja alfabetizaci{\'o}n; la traducci{\'o}n NLLB adaptada al quechua de Cusco (\texttt{quz}) es el n{\'u}cleo del aporte. Elaboraci{\'o}n propia.}
\label{fig:arq}
\end{figure}

\begin{algorithm}[H]
\caption{Traducci{\'o}n bidireccional por voz de un turno de di{\'a}logo}
\label{alg:pipeline}
\begin{algorithmic}[1]
\Require Audio de entrada $a$; idioma origen $L_s$; idioma destino $L_t$
\Ensure Texto traducido $t_t$ y audio sintetizado $a_t$
\State $t_s \gets \textsc{STT}_{\mathrm{MMS}}(a,\,L_s)$ \Comment{voz a texto}
\State $t_t \gets \textsc{MT}_{\mathrm{NLLB}}(t_s,\,L_s \rightarrow L_t)$ \Comment{traducci{\'o}n (\texttt{quz}$\leftrightarrow$es)}
\State Mostrar $t_t$ en pantalla \Comment{salida en texto}
\State $a_t \gets \textsc{TTS}_{\mathrm{MMS}}(t_t,\,L_t)$ \Comment{texto a voz}
\State Reproducir $a_t$ \Comment{salida en voz: inclusi{\'o}n de analfabetos}
\State \Return $(t_t,\,a_t)$
\end{algorithmic}
\end{algorithm}

\subsection{Hip{\'o}tesis}
\textbf{General:} un sistema de traducci{\'o}n bidireccional Quechua--Espa{\~n}ol por voz,
basado en modelos de IA preentrenados, mejora significativamente la comunicaci{\'o}n
m{\'e}dico-paciente en consultas hospitalarias de Cusco frente a la ausencia de apoyo tecnol{\'o}gico.
\begin{itemize}
  \item \textbf{H1:} un corpus cl{\'i}nico Quechua--Espa{\~n}ol validado proporciona la base l{\'e}xica necesaria para el ajuste de los modelos.
  \item \textbf{H2:} el ajuste (fine-tuning) de los modelos preentrenados reduce el WER y mejora el BLEU/chrF frente a los modelos sin ajustar.
  \item \textbf{H3:} la salida texto+voz mejora la comprensi{\'o}n de las indicaciones frente a solo texto, sobre todo con baja alfabetizaci{\'o}n.
  \item \textbf{H4:} el uso del sistema mejora significativamente la comprensi{\'o}n mutua percibida frente a la ausencia de apoyo tecnol{\'o}gico.
\end{itemize}

\subsection{Delimitaci{\'o}n}
\textbf{Espacial:} un centro de salud de referencia en Cusco. \textbf{Temporal:} octubre 2026
a enero 2027. \textbf{Tem{\'a}tica y {\'e}tica:} vocabulario acotado de consulta general, con
di{\'a}logos simulados (no datos reales de pacientes), en cumplimiento del atributo de {\'e}tica
del curso.

\subsection{Organizaci{\'o}n del documento}
El presente plan de tesis se organiza en tres cap{\'i}tulos y secciones complementarias. El
\textbf{Cap{\'i}tulo~I} expone la realidad problem{\'a}tica, la formulaci{\'o}n del problema, los
objetivos, la arquitectura propuesta, la justificaci{\'o}n, las hip{\'o}tesis y la delimitaci{\'o}n
del estudio. El \textbf{Cap{\'i}tulo~II} sintetiza la revisi{\'o}n de literatura que sustenta el
aporte ---su desarrollo completo, con el protocolo y el an{\'a}lisis de los estudios, se presenta
en el documento \texttt{01\_Review}---. El \textbf{Cap{\'i}tulo~III} detalla la metodolog{\'i}a:
enfoque, dise{\~n}o, variables, m{\'e}tricas, poblaci{\'o}n y muestra, procedimiento, an{\'a}lisis
de datos y aspectos {\'e}ticos. Finalmente, la \textbf{matriz de consistencia}, el
\textbf{cronograma} y los \textbf{anexos} respaldan la coherencia y la ejecuci{\'o}n del proyecto.

% ============================================================
\section{Cap{\'i}tulo II: Revisi{\'o}n de literatura}
El desarrollo extenso de la revisi{\'o}n sistem{\'a}tica ---protocolo, b{\'u}squeda y an{\'a}lisis
de los estudios--- se presenta en el documento \texttt{01\_Review}. Aqu{\'i} se sintetiza solo lo
esencial para situar el aporte, sin repetir lo ya expuesto en el Cap{\'i}tulo~I.

La literatura confirma dos hechos y un vac{\'i}o. Primero, la barrera ling{\"u}{\'i}stica
deteriora la atenci{\'o}n de pacientes quechuahablantes, documentada en hospitales de Cusco
\parencite{begazo2026} y asociada a la calidad de atenci{\'o}n percibida \parencite{cuaila2022}.
Segundo, ya existe tecnolog{\'i}a base de voz y traducci{\'o}n para lenguas de bajo recurso
---MMS \parencite{pratap2024mms}, NLLB \parencite{costajussa2022nllb} y antecedentes de
traducci{\'o}n neural del quechua \parencite{ortega2020quechua,mager2021americasnlp}, junto con
corpus como \emph{Huqariq} \parencite{zevallos2022huqariq}. El vac{\'i}o es su intersecci{\'o}n:
no existe un corpus cl{\'i}nico Quechua--Espa{\~n}ol ni una evaluaci{\'o}n emp{\'i}rica de estos
modelos en el dominio m{\'e}dico, y NLLB a{\'u}n no cubre la variante de Cusco (\texttt{quz}).
Cerrar esa intersecci{\'o}n es el aporte de este proyecto.

% ============================================================
\section{Cap{\'i}tulo III: Metodolog{\'i}a}

\subsection{Enfoque, tipo y nivel}
La investigaci{\'o}n adopta un enfoque \textbf{cuantitativo con complemento cualitativo}
(mixto), de tipo \textbf{aplicado} ---pues busca resolver un problema concreto mediante un
prototipo--- y de nivel \textbf{explicativo}, ya que contrasta el efecto del sistema frente a la
ausencia de apoyo tecnol{\'o}gico.

\subsection{Dise{\~n}o de investigaci{\'o}n}
Se emplea un dise{\~n}o \textbf{experimental de tipo pre-experimental / cuasi-experimental},
con mediciones de desempe{\~n}o del sistema y comparaci{\'o}n de modalidades de salida. La
unidad de an{\'a}lisis es el \emph{turno de traducci{\'o}n} dentro de una consulta simulada.

\subsection{Variables}
\begin{table}[H]
\centering
\caption{Operacionalizaci{\'o}n de variables.}
\label{tab:variables}
\small
\begin{tabularx}{\linewidth}{@{}l X l@{}}
\toprule
\textbf{Variable} & \textbf{Definici{\'o}n / dimensi{\'o}n} & \textbf{Indicador} \\
\midrule
\multirow{2}{*}{Independiente} & Uso del sistema de traducci{\'o}n IA & Con / sin sistema \\
                               & Modalidad de salida & Texto / texto+voz \\
\midrule
\multirow{3}{*}{Dependiente}   & Precisi{\'o}n de transcripci{\'o}n & WER \\
                               & Precisi{\'o}n de traducci{\'o}n & BLEU / chrF \\
                               & Comprensi{\'o}n y percepci{\'o}n & \% aciertos, escala Likert, MOS \\
\midrule
Interviniente                  & Nivel de alfabetizaci{\'o}n del participante & Alto / bajo \\
\bottomrule
\end{tabularx}
\end{table}

\subsection{M{\'e}tricas de evaluaci{\'o}n}
La precisi{\'o}n del reconocimiento de voz se mide con la \emph{tasa de error de palabra} (WER),
que relaciona las sustituciones $S$, inserciones $I$ y eliminaciones $D$ con el total de
palabras de referencia $N$:
\begin{equation}
\mathrm{WER} = \frac{S + D + I}{N}.
\end{equation}
\regecuacion{Tasa de error de palabra (WER)}
La calidad de la traducci{\'o}n se eval{\'u}a con BLEU \parencite{papineni2002bleu} y con chrF
\parencite{popovic2015chrf}; esta {\'u}ltima es adecuada para lenguas morfol{\'o}gicamente ricas
porque opera sobre $n$-gramas de caracteres:
\begin{equation}
\mathrm{chrF}_{\beta} = (1 + \beta^{2})\,\frac{\mathrm{chrP}\cdot\mathrm{chrR}}{\beta^{2}\,\mathrm{chrP} + \mathrm{chrR}},
\end{equation}
\regecuacion{Puntaje chrF ($F$-score de $n$-gramas de caracteres)}
donde $\mathrm{chrP}$ y $\mathrm{chrR}$ son la precisi{\'o}n y la exhaustividad de caracteres y
$\beta$ pondera su importancia relativa. La calidad de la voz sintetizada se valora con el
\emph{Mean Opinion Score} (MOS).

\subsection{Poblaci{\'o}n y muestra}
La validaci{\'o}n se realiza con \textbf{participantes voluntarios} (estudiantes de ciencias de
la salud y/o personal asistencial bilingüe) que \emph{act{\'u}an} roles de personal y de
paciente en consultas simuladas. \textbf{No se utilizan pacientes reales ni datos cl{\'i}nicos
reales.} El muestreo es \textbf{no probabil{\'i}stico por conveniencia}; se estima un m{\'i}nimo
de 10--15 participantes y un conjunto de 30--50 frases-gu{\'i}on por escenario.

\subsection{T{\'e}cnicas e instrumentos}
\begin{itemize}
  \item \textbf{Corpus:} lista de vocabulario y frases de triaje/s{\'i}ntomas/indicaciones,
        construida y revisada con personal de salud (fichas de validaci{\'o}n de contenido).
  \item \textbf{Registro de desempe{\~n}o:} \emph{logs} autom{\'a}ticos del sistema (WER, BLEU,
        latencia por turno).
  \item \textbf{Encuesta de percepci{\'o}n:} cuestionario tipo Likert (utilidad, confianza,
        comprensi{\'o}n) al personal y a los pacientes-actores.
  \item \textbf{Prueba de comprensi{\'o}n:} tarea de recuperaci{\'o}n de la indicaci{\'o}n
        (\% de aciertos) para contrastar H2 y H3.
\end{itemize}

\subsection{Procedimiento}
\begin{enumerate}
  \item Construcci{\'o}n y validaci{\'o}n del corpus m{\'e}dico Quechua--Espa{\~n}ol.
  \item Ajuste (fine-tuning) de MMS (reconocimiento y s{\'i}ntesis de voz) y, como aporte central, de NLLB para el par quechua de Cusco (\texttt{quz})--Espa{\~n}ol.
  \item Integraci{\'o}n del pipeline y de la salida diferenciada (Figura~\ref{fig:arq}).
  \item Ejecuci{\'o}n de consultas simuladas bajo las condiciones experimentales.
  \item Registro de m{\'e}tricas y aplicaci{\'o}n de instrumentos de percepci{\'o}n.
\end{enumerate}

\subsection{An{\'a}lisis de datos}
Las m{\'e}tricas t{\'e}cnicas (WER, BLEU, MOS, latencia) se resumen con estad{\'i}stica
descriptiva. Para contrastar hip{\'o}tesis (p. ej. modelo ajustado vs. sin ajustar, o
texto vs. texto+voz) se aplican pruebas de comparaci{\'o}n de medias/proporciones (t de Student
o Wilcoxon seg{\'u}n normalidad), con un nivel de significancia $\alpha = 0{,}05$.

\subsection{Aspectos {\'e}ticos}
El estudio usa {\'u}nicamente di{\'a}logos \textbf{simulados/actuados}, sin informaci{\'o}n
cl{\'i}nica real, evitando la necesidad de aprobaci{\'o}n de un comit{\'e} de {\'e}tica
hospitalario para datos de pacientes. Se obtiene consentimiento informado de los participantes
voluntarios y se anonimizan las grabaciones, en cumplimiento del atributo de {\'e}tica del curso.

% ------------------- MATRIZ DE CONSISTENCIA -------------------
\section{Matriz de consistencia}
\begin{table}[H]
\centering
\caption{Matriz de consistencia del proyecto.}
\label{tab:matriz}
\scriptsize
\renewcommand{\arraystretch}{1.3}
\begin{tabularx}{\linewidth}{@{}p{3.1cm} X X X@{}}
\toprule
\textbf{Problema} & \textbf{Objetivo} & \textbf{Hip{\'o}tesis} & \textbf{Variable/Indicador} \\
\midrule
General: {\textquestiondown}De qu{\'e} manera el sistema influye en la comunicaci{\'o}n m{\'e}dico-paciente? &
Desarrollar y evaluar el sistema de traducci{\'o}n bidireccional por voz. &
El sistema mejora significativamente la comunicaci{\'o}n m{\'e}dico-paciente. &
Comprensi{\'o}n (\% aciertos, Likert). \\
\midrule
E1: {\textquestiondown}Qu{\'e} corpus cl{\'i}nico se requiere? &
Construir y validar el corpus Quechua--Espa{\~n}ol. &
H1: el corpus validado aporta la base l{\'e}xica para el ajuste. &
Cobertura l{\'e}xica, validaci{\'o}n de contenido. \\
\midrule
E2: {\textquestiondown}El ajuste mejora la precisi{\'o}n? &
Adaptar NLLB al quechua de Cusco (\texttt{quz}) y MMS al dominio cl{\'i}nico. &
H2: el ajuste reduce el WER y mejora el BLEU/chrF. &
WER, BLEU/chrF. \\
\midrule
E3: {\textquestiondown}Qu{\'e} modalidad de salida es m{\'a}s comprensible? &
Integrar la salida por voz y texto bidireccional para ambos interlocutores. &
H3: texto+voz mejora la comprensi{\'o}n. &
\% aciertos, MOS. \\
\midrule
E4: {\textquestiondown}Qu{\'e} utilidad/confianza perciben los usuarios? &
Validar en consulta simulada y medir percepci{\'o}n. &
H4: mejora la comprensi{\'o}n mutua percibida. &
Escala Likert, tiempo. \\
\bottomrule
\end{tabularx}
\end{table}

% ============================================================
\section{Cronograma}
\begin{table}[H]
\centering
\caption{Cronograma tentativo del proyecto (2026--2027).}
\label{tab:cronograma}
\small
\begin{tabularx}{\linewidth}{@{}X c c c c@{}}
\toprule
\textbf{Actividad} & \textbf{Oct} & \textbf{Nov} & \textbf{Dic} & \textbf{Ene} \\
\midrule
Revisi{\'o}n de literatura      & $\bullet$ & $\bullet$ &           &           \\
Construcci{\'o}n del corpus     &           & $\bullet$ & $\bullet$ &           \\
Adaptaci{\'o}n de modelos       &           & $\bullet$ & $\bullet$ & $\bullet$ \\
Validaci{\'o}n simulada         &           &           & $\bullet$ & $\bullet$ \\
Redacci{\'o}n final             &           &           &           & $\bullet$ \\
\bottomrule
\end{tabularx}
\end{table}

% ------------------- BIBLIOGRAFIA -------------------
\newpage
\printbibliography[heading=bibintoc,title={Referencias}]

% ------------------- ANEXOS -------------------
\newpage
\anexo{Modelo de consentimiento informado para participantes voluntarios}
\textit{[Pendiente de elaboraci{\'o}n: formato de consentimiento conforme al atributo de
{\'e}tica del curso.]}

\anexo{Gui{\'o}n de consulta simulada (ejemplo de frases de triaje)}
\textit{[Pendiente de elaboraci{\'o}n: listado de 30--50 frases-gui{\'o}n por escenario.]}

\end{document}
